%!TEX root = ../main.tex

\section{Процедура Колмана--Нолла}

Подставляя уравнение баланса микросил и микронапряжений \eqref{eq:microforces} в закон сохранения внутренней энергии \eqref{eq:energy_conservation}, получаем
\[
    \dt{e} = -\Div \vq - \Div[\vE, \vect{H}] + \vxi \cdot \nabla \dt{\phi} - \pi \dt{\phi} - \Div(\kappa \rho \vj) + r.
\]
Выражая $\dt{e}$ из формулы \eqref{eq:free_energy_density}, получаем уравнение баланса свободной энергии в виде
\begin{equation}
    \dt{\psi} + \dt{\theta} \eta + \theta \dt{\eta} + \dt{\vE} \cdot \vD + \vE \cdot \dt{\vD} = -\Div \vq - \Div[\vE, \vect{H}] + \vxi \cdot \nabla \dt{\phi} - \pi \dt{\phi} - \Div(\kappa \rho \vj) + r.
    \label{eq:free_energy_conservation}
\end{equation}

Домножим энтропийное неравенство \eqref{eq:entropy_inequality} на $\theta$ и вычтем из уравнения \eqref{eq:free_energy_conservation}, получим
\[
    \dt{\psi} + \dt{\theta} \eta + \dt{\vE} \cdot \vD + \vE \cdot \dt{\vD} \leqslant -\Div \vq + \theta \Div \left( \cfrac{\vq}{\theta} \right) + \vE \cdot (\vj + \dt{\vD}) + \vxi \cdot \nabla \dt{\phi} - \pi \dt{\phi} - \Div(\kappa \rho \vj)
\]
и, после сокращения,
\begin{equation}
    \dt{\psi} + \dt{\theta} \eta + \dt{\vE} \cdot \vD + \cfrac{\vq}{\theta} \cdot \nabla \theta - \vE \cdot \vj - \vxi \cdot \nabla \dt{\phi} + \pi \dt{\phi} + \Div(\kappa \rho \vj) \leqslant 0.
    \label{eq:free_energy_inequality}
\end{equation}

Постулированная ранее функциональная зависимость $\psi = \hat{\psi}(\theta, \vE, \phi, \nabla \phi, \rho)$ влечет, по формуле дифференцирования сложной функции,
\[
    \dt{\psi} = \hat{\psi}_\theta \dt{\theta} + \hat{\psi}_\vE \cdot \dt{\vE} + \hat{\psi}_\phi \dt{\phi} + \hat{\psi}_{\nabla \phi} \cdot \nabla \dt{\phi} + \hat{\psi}_\rho \dt{\rho}.
\]
Подставляя это выражение в неравенство \eqref{eq:free_energy_inequality}, имеем
\[
    \dt{\theta} (\eta + \hat{\psi}_\theta) + \dt{\vE} \cdot (\vD + \hat{\psi}_\vE) + \cfrac{\vq}{\theta} \cdot \nabla \theta - \vE \cdot \vj + \nabla \dt{\phi} \cdot (-\vxi + \hat{\psi}_{\nabla \phi}) + \dt{\phi}(\pi + \hat{\psi}_\phi) + \Div(\kappa \rho \vj) + \hat{\psi}_\rho \dt{\rho} \leqslant 0.
\]
Положим
\[
    \hat{\psi}_\theta = -\eta, \quad \hat{\psi}_\vE = -\vD, \quad \hat{\psi}_{\nabla \phi} = \vxi,
\]
что обращает в ноль соответствующие слагаемые в неравенстве. Пусть также
\[
    \hat{\psi}_\rho = \kappa \rho.
\]
Тогда, учитывая, что
\[
    \Div(\kappa \rho \vj) + \kappa \rho \dt{\rho} = \Div(\kappa \rho \vj) - \kappa \rho \Div \vj = \nabla(\kappa \rho) \cdot \vj,
\]
получим
\[
    \cfrac{\vq}{\theta} \cdot \nabla \theta + \dt{\phi}(\pi + \hat{\psi}_\phi) - [\vE - \nabla(\kappa \rho)] \cdot \vj \leqslant 0.
\]
Далее положим
\[
    \hat{\psi}_\phi = -\pi - \beta \dt{\phi}, \quad \vq = -k \nabla \theta, \quad \vj = \sigma [\vE - \nabla (\kappa \rho)],
\]
где $\beta, k, \sigma, \kappa > 0$ и, вообще говоря,
\[
    \beta = \beta(\theta, \phi), \quad k = k(\theta, \phi), \quad \sigma = \sigma(\theta, \phi), \quad \kappa = \kappa(\theta, \phi).
\]
Неравенство для свободной энергии, таким образом, приобретает вид
\[
    -\cfrac{k}{\theta} \nabla \theta \cdot \nabla \theta - \beta \dt{\phi}^2 - \sigma [\vE - \nabla(\kappa \rho)] \cdot [\vE - \nabla(\kappa \rho)] \leqslant 0,
\]
оказываясь верным при любых значениях свободных переменных.

Учитывая полученные выражения для частных производных $\hat{\psi}$, можно записать
\[
    \dt{\psi} = -\dt{\theta} \eta - \dt{\vE} \cdot \vD + \vxi \cdot \nabla \dt{\phi} - \dt{\phi} (\pi + \beta \dt{\phi}) - \kappa \rho \Div \vj.
\]
Отыщем выражение для $\dt{\theta}$. Подставим соотношение выше в уравнение \eqref{eq:free_energy_conservation} баланса свободной энергии:
\begin{gather*}
    \theta \dt{\eta} + \vE \cdot \dt{\vD} - \beta \dt{\phi}^2 - \kappa \rho \Div \vj = -\Div \vq + \vE \cdot (\vj + \dt{\vD}) - \Div(\kappa \rho \vj) + r, \\
    \theta \dt{\eta} = -\Div \vq + \beta \dt{\phi}^2 + [\vE - \nabla(\kappa \rho)] \cdot \vj + r.
\end{gather*}
Так как $\eta = -\hat{\psi}_\theta$, то
\[
    \theta \dt{\eta} = -\theta \left[ \hat{\psi}_{\theta \theta} \dt{\theta} + \hat{\psi}_{\theta \vE} \cdot \dt{\vE} + \hat{\psi}_{\theta \phi} \dt{\phi} + \hat{\psi}_{\theta \nabla \phi} \cdot \nabla \dt{\phi} + \hat{\psi}_{\theta \rho} \dt{\rho} \right].
\]
Величину $-\theta \hat{\psi}_{\theta \theta}$ будем обозначать $c$ и называть теплоемкостью. Можно выразить
\begin{equation}
    c \dt{\theta} = -\Div \vq + \beta \dt{\phi}^2 + [\vE - \nabla(\kappa \rho)] \cdot \vj + r + \theta \left[ \hat{\psi}_{\theta \vE} \cdot \dt{\vE} + \hat{\psi}_{\theta \phi} \dt{\phi} + \hat{\psi}_{\theta \nabla \phi} \cdot \nabla \dt{\phi} + \hat{\psi}_{\theta \rho} \dt{\rho} \right],
    \label{eq:temperature}
\end{equation}
что в описываемой модели является уравнением динамики температуры $\theta$.

Отметим, что уравнение \eqref{eq:temperature} гарантирует (при нулевом $r$) $\theta > 0$. Добавление диффузии заряда в описанном виде не нарушило этого свойства: третье слагаемое неотрицательно.

Уравнение динамики фазового поля $\phi$ получим из баланса \eqref{eq:microforces} микросил и микронапряжений. Подставляя $\vxi = \hat{\psi}_{\nabla \phi}$ и $\pi = -\hat{\psi}_\phi - \beta \dt{\phi}$, имеем
\begin{equation}
    \beta \dt{\phi} = \Div \hat{\psi}_{\nabla \phi} - \hat{\psi}_\phi + \gamma.
    \label{eq:phi}
\end{equation}


\section{Дальнейшее уточнение соотношений}

Примем следующую форму зависимости свободной энергии $\psi$ от параметров состояния среды:
\[
    \hat{\psi}(\theta, \vE, \phi, \nabla \phi, \rho) = \psi^{(2)}(\theta, \vE, \phi, \rho) + \cfrac{\lambda}{2} \nabla \phi \cdot \nabla \phi.
\]
Пусть $\lambda = \Const$, $\kappa = \Const$. Пусть также $\vD = \epsilon(\theta, \phi) \vE$. Тогда из $\psi^{(2)}_\vE = \hat{\psi}_\vE = -\vD = -\epsilon \vE$ следует
\[
    \psi^{(2)} = -\half \epsilon(\theta, \phi) \vE \cdot \vE + \psi^{(1)}(\theta, \phi, \rho).
\]
Аналогично, $\psi^{(1)}_\rho = \hat{\psi}_\rho = \kappa \rho$ влечет
\[
    \psi^{(1)} = \cfrac{\kappa}{2} \rho^2 + \psi^{(0)}(\theta, \phi).
\]
Таким образом, формула плотности свободной энергии имеет вид
\begin{equation}
    \hat{\psi}(\theta, \vE, \phi, \nabla \phi, \rho) = -\half \epsilon(\theta, \phi) \vE \cdot \vE + \psi^{(0)}(\theta, \phi) + \cfrac{\kappa}{2} \rho^2 + \cfrac{\lambda}{2} \nabla \phi \cdot \nabla \phi.
\end{equation}
При этом определяющие соотношения модели принимают вид
\begin{align*}
    \eta &= \half \epsilon_\theta(\theta, \phi) \vE \cdot \vE - \psi^{(0)}_\theta(\theta, \phi), \\
    \vD &= \epsilon(\theta, \phi) \vE, \\
    \vxi &= \lambda \nabla \phi, \\
    \pi &= \half \epsilon_\phi(\theta, \phi) \vE \cdot \vE - \psi^{(0)}_\phi(\theta, \phi) - \beta(\theta, \phi) \dt{\phi}, \\
    \vq &= -k(\theta, \phi) \nabla \theta, \\
    \vj &= \sigma(\theta, \phi) (\vE - \kappa \nabla \rho).
\end{align*}


\section{Определение функций свойств среды}

Для замыкания уравений модели остается определить функции $\epsilon$, $\sigma$, $\beta$, $k$ и $\psi^{(0)}$ переменных $\theta$ и $\phi$, выражающие зависимость свойств среды от температуры и фазового поля, то есть связь этих свойств с фазой вещества.

Для функций свойств среды $\epsilon$, $\sigma$ в данной точке пространства постулируем зависимость исключительно от фазового поля $\phi$. Считая $\epsilon(\vx, \phi = 1) = \epsilon_0(\vx)$ и $\sigma_0(\vx, \phi = 1) = \sigma_0(\vx)$ известными, а также полагая $\epsilon(\phi = 0) \gg \epsilon(\phi = 1)$, $\sigma(\phi = 0) \gg \sigma(\phi = 1)$, примем
\begin{align*}
    \epsilon[\phi] &= \cfrac{\epsilon_0(\vx)}{g(\phi) + \delta_\epsilon}, \\
    \sigma[\phi] &= \cfrac{\sigma_0(\vx)}{g(\phi) + \delta_\sigma},
\end{align*}
где $g(\phi)$ -- интерполирующая функция, гладкая и возрастающая, соединяющая значения $g(0) = 0$ и $g(1) = 1$; $\delta_\epsilon, \delta_\sigma \ll 1$, константы.

Положим $\beta(\theta, \phi) = 1 / m = \Const$, где $m$ задает скорость изменения $\phi$ под действием единичной обобщенной силы и называется подвижностью.

Пусть коэффициент теплопроводности $k(\theta, \phi) = k = \Const$.

Примем следующий вид функции $\psi^{(0)}$:
\[
    \psi^{(0)}(\theta, \phi) = \psi^T(\theta, \phi) + \Gamma_V [1 - f(\phi)].
\]
Здесь первое слагаемое $\psi^T$ есть энергия, связанная с температурой среды. Второе слагаемое -- энергия, затрачиваемая на разрушение среды, то есть на фазовый переход от $\phi = 1$ к $\phi = 0$; $f(\phi)$ -- интерполирующая функция, гладкая и возрастающая, соединяющая значения $f(0) = 0$, $f(1) = 1$; $\Gamma_V$ имеет смысл энергии, затрачиваемой на фазовый переход единицы объема вещества.

Для вывода формулы $\phi^T$ будем считать, что теплоемкость среды $c = -\hat{\psi}_{\theta \theta} = -\psi^T_{\theta \theta} \hm = c(\phi)$, то есть не зависит от температуры. Тогда, интегрируя, получим
\[
    \psi^T = c(\phi) \theta (1 - \ln \theta) + \widetilde{C}(\phi),
\]
где $\widetilde{C}$ -- произвольная функция $\phi$. Для энтропии $\eta = -\psi^T_\theta$ следует формула
\[
    \eta = c(\phi) \ln \theta + \widetilde{C}(\phi).
\]
Имея опорные значения энтропии $\eta_0(\phi)$ для $\phi \in [0, 1]$ при заданной $\theta_0$, можно выразить
\begin{align*}
    \eta(\theta, \phi) &= c(\phi) \ln \cfrac{\theta}{\theta_0} + \eta_0(\phi), \\
    \psi^T(\theta, \phi) &= \theta \left[ c(\phi) - c(\phi) \ln \cfrac{\theta}{\theta_0} - \eta_0(\phi) \right].
\end{align*}

Отметим, что величина $l = \sqrt{\lambda / \Gamma_V}$ есть масштабный коэффициент характерной толщины диффузной границы. Переобозначим $\Gamma_V = \Gamma / l^2$, $\lambda = \Gamma$; здесь параметр $\Gamma$ имеет смысл характерной энергии образования электрического пробоя единичной длины (подразумевается <<площадь>> сечения $\approx l^2$).